\documentclass[UTF8]{ctexart}
\usepackage{xeCJK}
\usepackage{geometry}
\geometry{left=2cm,right=2cm,top=2.5cm,bottom=2.5cm}
\setlength{\headheight}{12.7pt}

\usepackage{titlesec}
\titleformat{\section}
  {\normalfont\large\bfseries}
  {\thesection}
  {1em}
  {}
\titlespacing*{\section}
  {0pt}
  {3.5ex plus 1ex minus .2ex}
  {2.3ex plus .2ex}

\usepackage{amsmath}
\usepackage{amssymb}
\usepackage{amsthm}
\usepackage{enumitem}
\setlist[enumerate]{itemsep=0pt,topsep=0pt,parsep=0pt,leftmargin=0pt,itemindent=2.5em}
\setlist[itemize]{itemsep=0pt,topsep=0pt,parsep=0pt,leftmargin=0pt,itemindent=2.5em}
\usepackage{fancyhdr}
\pagestyle{fancy}
\fancyhf{}
\usepackage{graphicx}
\usepackage{hyperref}
\usepackage{indentfirst}
\usepackage{lmodern}

\begin{document}
\setlength{\abovedisplayskip}{1pt}
\setlength{\belowdisplayskip}{1pt}

\section{滤子基}

% 旅游好累, 难得休息一天. 可能写的不够好, 日后再做改进.
% 2026.05.05 于许昌

\hypertarget{sec:定义1.1}{}
\noindent\textbf{定义1.1 (滤子基)}

给定集合$X$和它的非空的子集族$B\subset\mathcal{P}(X)$. 如果:
\begin{enumerate}
    \item $\varnothing\notin B$,
    \item 对任意的$b_1,b_2\in B$, 存在$b_3\in B$使得
    $b_3\subset b_1\cap b_2$,
\end{enumerate}
就称$B$是$X$上的一个\textbf{滤子基}.

\vspace{10pt}

滤子基通过向上扩张操作就可以得到滤子.

\vspace{10pt}

\hypertarget{sec:定义1.2}{}
\noindent\textbf{定义1.2 (生成滤子)}

任给集合$X$和$X$上的滤子基$B$,
\[
    \langle B\rangle=\{a\subset X\mid\exists\,b\in B:b\subset a\}
\]
是$X$上包含$B$的最小滤子, 称作由$B$生成的滤子.
也称$B$是$\langle B\rangle$的一个基.

\noindent{\kaishu 验证.}

显然$\langle B\rangle$包含$B$. 首先证明$\langle B\rangle$是滤子.

\begin{enumerate}
    \item 显然$\varnothing\notin\langle B\rangle,\,X\in\langle B\rangle$.
    \item 如果$a_1\in\langle B\rangle$且$a_1\subset a_2\subset X$,
    那么存在$b\in B$使得
    \[
        b\subset a_1\quad\implies\quad b\subset a_2,
    \]
    从而$a_2\in\langle B\rangle$.
    \item 如果$a_1,a_2\in\langle B\rangle$, 那么存在$b_1,b_2\in B$使得
    \[
        b_1\subset a_1\quad\land\quad b_2\subset a_2,
    \]
    此时存在$b_3\in B$使得$b_3\subset b_1\cap b_2\subset a_1\cap a_2$,
    所以$a_1\cap a_2\in\langle B\rangle$.
\end{enumerate}

接下来, 如果滤子$F$包含$B$, 那么对任意的$a\in\langle B\rangle$,
存在$b\in B$使得$b\subset a$, 此时
\[
    b\in F\quad\implies\quad a\in F,
\]
所以$\langle B\rangle\subset F$, 即$\langle B\rangle$是$X$上包含$B$的最小滤子.





\newpage

\section{滤子子基}

\hypertarget{sec:定义2.1}{}
\noindent\textbf{定义2.1 (滤子子基)}

给定集合$X$和它的非空的子集族$S\subset\mathcal{P}(X)$.

如果$S$满足``有限交性质'', 即
$S$的非空有限子集的交集都非空, 那么就称$S$是$X$上的一个\textbf{滤子子基}.

\vspace{10pt}

滤子子基通过有限交扩张操作可以得到滤子基, 从而可以生成滤子.

\vspace{10pt}

\hypertarget{sec:定义2.2}{}
\noindent\textbf{定义2.2 (生成滤子)}

任给集合$X$和$X$上的滤子子基$S$,
\vspace{3pt}
\[
    B=\left\{\,\bigcap R\,\,\middle\vert\,\,
    R\text{\,是\,}S\text{\,的非空有限子集\,}\right\}\\[3pt]
\]
是$X$上包含$S$的滤子基.

进一步, $\langle B\rangle$是$X$上包含$S$的最小滤子, 也记作$\langle S\rangle$,
称作由$S$生成的滤子, 称$S$是其子基.

\noindent{\kaishu 验证.}

显然$B$包含$S$. 首先证明$B$是滤子基.

\begin{enumerate}
    \item 依定义, $S$的非空有限子集的交集非空, 所以$\varnothing\notin B$.
    \item 对任意的$b_1,b_2\in B$, 存在$S$的非空有限子集$R_1,R_2$使得
    $\displaystyle b_1=\bigcap R_1,\,b_2=\bigcap R_2$, 从而
    \vspace{3pt}
    \[
        b_1\cap b_2=\bigcap\,(R_1\cup R_2)\in B.
    \]
\end{enumerate}

% Maybe we won't see the sunrise
% Maybe we won't see the moon shine
% Even if the whole world is just you and I
% It's ok girl everything will be fine

\vspace{3pt}
接下来, 如果滤子$F$包含$S$, 那么对任意的$b\in B$, 存在$S$的非空有限子集$R$使得
\[
    b=\bigcap R\in F,
\]
所以$B\subset F$, 从而$\langle B\rangle\subset F$,
即$\langle B\rangle$是$X$上包含$S$的最小滤子. \hfill $\qedsymbol$

\end{document}
